\maketableofcontents

\chapter*{\textbf{ВВЕДЕНИЕ}}
\addcontentsline{toc}{chapter}{\textbf{ВВЕДЕНИЕ}}

В условиях активного развития сетевых технологий и роста объёма информации особое значение приобретают системы, обеспечивающие эффективное взаимодействие пользователей в режиме реального времени. 
Одной из актуальных задач современной информатики является организация совместной работы над документами и заметками, позволяющая нескольким пользователям одновременно редактировать общий контент в распределённой сетевой среде.

Традиционные подходы к совместному редактированию, основанные на последовательной передаче изменений или централизованной синхронизации, сталкиваются с рядом проблем, связанных с задержками передачи данных, конфликтами версий и необходимостью постоянного соединения с сервером, что особенно критично в распределённых системах.
Эти недостатки особенно заметны при работе с распределёнными системами и в условиях ограниченных сетевых ресурсов.

Одним из перспективных направлений решения данной задачи является разработка и внедрение протоколов совместного редактирования, обеспечивающих корректную синхронизацию изменений между участниками без потери данных и с минимальными задержками. 
Такой протокол позволяет пользователям одновременно вносить изменения в заметки, автоматически разрешать конфликты и сохранять согласованное состояние документа.

Цель работы – разработать протокол для совместного редактирования заметок.

В ходе выполнения курсового проекта необходимо решить следующие основные задачи:

\begin{itemize}
	\item проанализировать предметную область и определить основные проблемы, возникающие при организации совместного редактирования заметок в распределённых системах;
	\item провести обзор существующих решений и протоколов синхронизации данных, выявить их преимущества и недостатки;
	\item определить требования к разрабатываемому протоколу, включая аспекты согласованности данных, устойчивости к конфликтам и производительности;
    \item разработать структуру и формат сообщений, обеспечивающих обмен изменениями между участниками редактирования;
    \item реализовать прототип протокола, обеспечивающего синхронизацию заметок между несколькими пользователями в реальном времени;
    \item провести тестирование и анализ корректности работы протокола при одновременном внесении изменений, оценить эффективность предложенного решения.
\end{itemize}